% Generated draft: requires mathematical and source review.
\documentclass[11pt]{article}
\usepackage[margin=25mm]{geometry}
\usepackage{fontspec}
\setmainfont{Latin Modern Roman}
\usepackage{amsmath,amssymb,mathtools,unicode-math}
\setmathfont{Latin Modern Math}
\usepackage{enumitem,needspace,xurl,hyperref}
\setlength{\parindent}{0pt}
\setlength{\parskip}{5pt}
\setlength{\emergencystretch}{3em}
\newcommand{\N}{\mathbb{N}}
\newcommand{\Z}{\mathbb{Z}}
\newcommand{\Q}{\mathbb{Q}}
\newcommand{\R}{\mathbb{R}}
\newcommand{\C}{\mathbb{C}}
\newcommand{\F}{\mathbb{F}}
\newcommand{\A}{\mathbb{A}}
\newcommand{\PP}{\mathbb{P}}
\newcommand{\E}{\mathbb{E}}
\newcommand{\eps}{\varepsilon}
\newcommand{\dd}{\,\mathrm{d}}
\begin{document}
\section*{UnsolvedMath ID: 20002079}
\textbf{Source problem code:} catalogue entry
\subsection{Definitions and mathematical statement}
Question 1.21 records Margaret Symington's program of using locally toric structure to count holomorphic curves, inspired by Mikhalkin's amoeba methods. The note explicitly says that the author lacks good notes on the series of questions. No precise invariant, curve class, constraints or correspondence theorem is recoverable from the text alone.

The aid reports affirmative pieces of this program in unnamed 2026 work on tropical formulas for disk potentials of almost-toric four-manifolds. Its exact hypotheses and citation are absent. Locally toric and almost-toric, allowing focus-focus nodes, must not be treated as identical definitions without qualification.

The aid's independent diagnostic uses the integral local system of fiber cycles in a regular torus fibration. Choose a basis in which monodromy about a multiplicity-\(k\) focus-focus node, \(k\geq1\), is
\[
 T_k=\begin{pmatrix}1&k\\0&1\end{pmatrix}.
\]
This matrix is specified on the chosen cycle local system; the affine tangent or dual-cotangent convention may use a different matrix. A loop of winding \(n\in\mathbb Z\) transports \(v=(a,b)\) to \(T_k^nv\), with endpoint defect
\[
 (T_k^n-I)v=nkb\,e_1,
\]
where \(e_1\) is the invariant direction. Without additional edges the coefficient closes exactly when \(nkb=0\); for nonzero winding this forces \(b=0\). Zero winding is not subject to that restriction. If invariant prongs are added, their total boundary contribution must cancel this defect. In the convention where their contribution is \(w e_1\), closure means
\[
 w=-nkb.
\]
Equivalently the compensating signed amount has magnitude and orientation prescribed by the defect; changing incoming/outgoing conventions changes the displayed sign. Ordinary chartwise \(\mathbb Z^2\)-balancing alone misses this local-system condition.

The companion action-angle statement is local. In a trivialization with angle coordinates of period one and
\(\omega=\sum_i dp_i\wedge d\theta_i\), sweeping the fiber cycle \(v\) along an affine path \(p(s)\) gives signed symplectic area
\[
 \int_C\omega=\langle v,p(1)-p(0)\rangle.
\]
Angles of period \(2\pi\), opposite symplectic sign or reversed cylinder orientation alter the corresponding factor or sign. A path crossing monodromy requires compatible local trivializations and the preceding boundary bookkeeping.

These are necessary geometric consistency filters for a proposed counting formalism. They neither prove existence of holomorphic lifts nor define a full enumerative invariant or recover Symington's missing detailed questions.
\subsection{Short English statement}
The source is an explicitly incomplete research-program note. The supplied aid gives a useful necessary condition: tropical cycle coefficients must close in their monodromy local system, with exact compensating defect at a focus-focus node. It also records local swept-cycle area, while leaving the full curve-counting problem unspecified.
\subsection{Sources}
\begin{itemize}
\item Local source record: catalogue entry, catalogue entry; canonical \texttt{aim-geometry-notes.json}, \texttt{notes[416]}, Question 1.21. Title: ``Local toric bookkeeping and a focus-focus monodromy obstruction.''
\item Michael Hutchings, with help from Yasha Eliashberg and John Etnyre, \emph{Holomorphic curves in contact geometry}, AIM workshop problem list, version 15 October 2003; supplied reference \url{https://aimath.org/WWN/contactgeom2/contactgeom2.pdf}. The question series is attributed to Margaret Symington, motivated by Mikhalkin. This draft does not access a PDF or HTML source.
\item Supplied attempt 1 gives the monodromy-defect and action-angle lemmas and mentions 2026 tropical disk-potential results without author, title or full bibliographic reference.
\item Registry dated 14 August 2026: status ``open,'' classification ``OPEN-TRIAGE,'' statement status ``exact''; attempt context only, research-run status partially solved, rigor ``proved,'' novelty confidence low. The machine-generated, non-peer-reviewed aid requires independent verification before citation.
\end{itemize}
\end{document}
